\chapter*{La solution en une page}
\addcontentsline{toc}{chapter}{La solution en une page}
\markboth{La solution en une page}{}

\paragraph{Le problème.} Un client écrit à sa banque en langage libre : \og Je veux acheter un
appartement \fg, \og J'ai perdu ma carte hier soir \fg, \og Pourquoi j'ai payé 8\eur{} de frais ? \fg.
L'assistant doit comprendre, poser les bonnes questions, tenir compte de la situation du client, et
proposer une piste --- avec une IA open source, locale, la plus sobre possible, et dont chaque
conseil se justifie.

\paragraph{Notre réponse : une cascade.} Chaque étage ne passe la main au suivant que s'il ne sait
pas faire lui-même. Le calcul le plus cher n'est appelé qu'en dernier recours, et c'est mesuré.

\begin{center}\small
\begin{tabularx}{\textwidth}{@{}l L{3.6cm} Y@{}}
\toprule
\textbf{Étage} & \textbf{Qui fait le travail} & \textbf{Pourquoi pas un LLM} \\
\midrule
1. Comprendre & Mots-clés, sinon petit encodeur (118 M de paramètres) et régression logistique &
  Les cas évidents coûtent moins d'une milliseconde, les autres quelques millisecondes ; tout se
  mesure sur le jeu de test officiel de Banking77. \\
2. Converser & Parcours décrits en YAML : quoi demander, dans quel ordre & Les questions à poser se
  déduisent du profil : on ne demande jamais ce qu'on sait déjà. \\
3. Conseiller & Règles publiques (plafonds, 35\,\% d'endettement) et modèle d'appétence &
  Une décision financière doit être reproductible et justifiable, pas générée. \\
4. Agir & Outils : opposition, virement guidé, transmission au conseiller & Une action ne part
  qu'après confirmation explicite du client. \\
5. Formuler & Le seul appel LLM du tour, à la fin : \f{ministral-3:3b} (3,0 Go) sous garde-fou ;
  gabarits pour les réponses simples & Le LLM met en forme des faits calculés ; il n'en invente
  aucun. \\
Escalade & \f{ministral-3:14b} (9,1 Go), seulement si l'étage 1 doute encore après une question &
  Le \og gros modèle à un instant précis \fg{} que CGI a explicitement encouragé. \\
\bottomrule
\end{tabularx}
\end{center}

\paragraph{Les quatre décisions qui font la différence.}
\begin{enumerate}
  \item \textbf{Comprendre sans LLM.} Des mots-clés pour les cas évidents, un petit classifieur
        pour le reste. Il est entraîné sur Banking77 traduit en français, plus 8 intentions que
        Banking77 ne couvre pas : épargne, crédit immobilier, crédit à la consommation, prêt
        étudiant, virement à un proche, jargon, insatisfaction, hors sujet.
  \item \textbf{Décider sans LLM.} Chaque proposition vient avec ses raisons \emph{et} avec la
        liste des produits écartés et pourquoi (\og pourquoi pas le LEP ? \fg).
  \item \textbf{Formuler avec le plus petit modèle qui suffit, une seule fois par tour.} On part
        d'un modèle qui parle bien (14B), on descend en taille, et on garde le plus petit qui tient
        encore le seuil de qualité.
  \item \textbf{Tout mesurer contre une version \og tout LLM \fg.} Même jeu de scénarios, trois
        configurations, et un tableau qualité / énergie / latence / taille au cœur du pitch.
\end{enumerate}

\paragraph{Ce qu'un organisateur nous a dit.} \og Pour gagner, c'est moins de LLM, plus d'algo, ou
des modèles d'IA plus spécialisés \fg, y compris pour le sujet 1 : des mots-clés et un classifieur
en entrée, un LLM appelé une seule fois, à la fin, et dont on choisit la taille en descendant.
C'est l'architecture ci-dessus.

\needspace{10\baselineskip}
\paragraph{Ce que la grille du jury y gagne.}
\begin{center}\small
\begin{tabularx}{\textwidth}{@{}l c Y@{}}
\toprule
\textbf{Critère} & \textbf{Poids} & \textbf{Notre preuve} \\
\midrule
Pertinence métier & 25\,\% & Les 6 situations des slides, plus l'appartement, le prêt étudiant et
  le jargon ; priorité n\textsuperscript{o}\,1 de La Banque Postale ; inclusion et accessibilité. \\
Qualité technique & 20\,\% & Précision sur le test officiel Banking77 (anglais et français) ; taux
  de réussite d'un jeu de 30 scénarios automatisés. \\
Frugalité & 20\,\% & Énergie par requête (CodeCarbon), latence, Go chargés, part des requêtes
  traitées sans LLM ; comparaison à la version \og tout LLM \fg. \\
Explicabilité & 15\,\% & Panneau \og pourquoi \fg{} : intention et confiance, règles appliquées,
  produits proposés et écartés, sources publiques. \\
Démonstration & 10\,\% & Six minutes en direct, avec les profils fournis par les slides ; vidéo de
  secours. \\
Présentation & 10\,\% & Trois voix, une par module ; un tableau de résultats, une démo, un
  message. \\
\bottomrule
\end{tabularx}
\end{center}

\paragraph{L'organisation.} Trois modules, un responsable et un suppléant chacun : \emph{Comprendre
et mesurer} (Younes), \emph{Converser et formuler} (Yahya), \emph{Conseiller et agir} (Sohayb).
Les contrats entre modules sont figés lundi soir ; deux intégrations mardi (12h30 et 18h) ; gel
fonctionnel mardi 18h ; mercredi matin consacré aux mesures finales, à la vidéo et aux répétitions.

\begin{cle}
\textbf{La phrase à retenir.} \emph{Notre assistant comprend souvent avec de simples mots-clés,
sinon avec un modèle de 118 millions de paramètres ; il décide avec des règles que le client peut
lire ; et il n'appelle un LLM qu'une fois, à la fin, pour mettre en mots ce qui a déjà été
calculé.}
\end{cle}
